\part{Hilbert概形}\label{fga3.iv}\thispagestyle{fancy}
\begin{original}{}\label{fga3.iv-original-citation}\skipsinglepar
\par{}A. Grothendieck. "Techniques de construction et théorèmes d'existence en géométrie algébrique, IV: Les schémas de Hilbert". \emph{Séminaire Bourbaki} \textbf{13} (1960–61), 第221讲, 1961年5月. \url{http://www.numdam.org/book-part/SB_1960-1961__6__249_0/}
\end{original}
\chapter{引言}\thispagestyle{fancy}\label{fga3.iv-introduction}
\par{}\hyperref[fga3.i]{[2, I]}和\hyperref[fga3.ii]{[2, II]}所述的方法, 大体上不依赖于对所讨论概形作任何射影性假设. 可惜它们至今尚未使我们解决\hyperref[fga3.ii]{[2, II]}所提出的存在性问题. 在本文和下一篇文章中, 我们将通过加上射影性假设来解决这些问题. 所用方法具有典型的射影性质, 实际上几乎不使用\hyperref[fga3.i]{[2, I]}和\hyperref[fga3.ii]{[2, II]}的结果. 这里将构造"Hilbert概形", 用来替代Chow坐标, 这在\hyperref[fga3.ii-c.2]{[2], II, C, §2}中已经提到. 在下一篇文章中, 将\hyperref[fga3.iii]{[2, III]}所建立的概形取商理论与Hilbert概形理论结合, 就能例如在相当一般的条件下构造Picard概形 (定义见\hyperref[fga3.ii-c.3]{[2], II, C, §3}).
\par{}总的说来, 现在我们已有一种大体令人满意的射影构造方法, 只是仍缺少\begingroup\renewcommand{\thefootnote}{*}\footnote{参见本讲末尾的增补。}\endgroup{}对"无不动点地"作用的射影群等群的取商理论 (参见\hyperref[fga3.iii-8]{[2], III, §8}). 解析几何的情形似乎还稍好一些 (若限于研究给定解析空间上的"射影"解析空间), 因为在解析空间中, 关于作用良好的群取商的困难已不再存在. 无论如何, 在代数几何和解析几何中, 都仍有待建立不需要任何射影性假设的构造方法.
\chapter{层的有界集合: 不变性质}\thispagestyle{fancy}\label{fga3.iv-1}
\par{}设\(k\)是域, \(X\)是\(k\)-预概形 (为简单起见, 假设它有限型). 每个扩张\(K/k\)给出一个\(K\)-预概形\(X_K=X\otimes _k K\). 若\(\mathscr {F}\)是\(X_K\)上的凝聚层, \(K'\)是\(K\)的一个扩张, 则\(\mathscr {F}\otimes _K K'=\mathscr {F}_{K'}\)是\(X_K\otimes _KK'=X_{K'}\)上的拟凝聚层. 因而若\(K\)、\(K'\)是\(k\)的任意扩张, \(\mathscr {F}\)是\(X_K\)上的拟凝聚层, \(\mathscr {F}'\)是\(X_{K'}\)上的拟凝聚层, 并且存在扩张\(K''/k\)和\(k\)-同态\(K\to  K''\)、\(K'\to  K''\), 使\(\mathscr {F}_{K''}\)与\(\mathscr {F}'_{K''}\)在\(X_{K''}\)上同构, 则称\(\mathscr {F}\)与\(\mathscr {F}'\)\emph{等价}. 这样定义了一个等价关系; 我们关心层在此关系下的等价类及这些等价类的集合. 注意, 若\(X_0\)在\(k\)上有限型, 则每个凝聚层的类均可由某个\(X_K\)上的凝聚层定义, 其中\(K\)是\(k\)的一个\emph{有限型}扩张. 因此, 在定义凝聚层的类时, 可以只考虑\(k\)的\emph{代数闭}扩张, 甚至可以只考虑\(k\)的一个固定的、具有无限超越次数的代数闭扩张\(\Omega \); 此时\(X_\Omega \)上两个凝聚层\(\mathscr {F}\)与\(\mathscr {F}'\)等价, 当且仅当存在\(\Omega \)的一个\(K\)-自同构\(\sigma \), 使\(\mathscr {F}\otimes _K(\Omega ,\sigma )\)同构于\(\mathscr {F}'\). 注意第一种定义下的凝聚层的类与第二种定义下的类之间有一个双射对应.
\par{}设\(E\)和\(E'\)是\(X\)上的两组凝聚层的类. 考虑所有形如\(\mathscr {F}\otimes \mathscr {F}'\)的层的类, 其中\(\mathscr {F}\)和\(\mathscr {F}'\)是\emph{同一个}的\(X_K\)上的凝聚层, 且\(\mathscr {F}\)的类属于\(E\), \(\mathscr {F}'\)的类属于\(E'\). 这样定义一个凝聚层的类的集合, 记作\(E\otimes  E'\). 同样可定义\(\mathscr {T}\kern  -2.75pt\mathit {or}_i(E,E')\)等. 一般地, 若一个运算\(\mathscr {U}\)把同一个\(X_K\)上的\(n\)个凝聚层组成的序列\(\mathscr {F}_1,\ldots ,\mathscr {F}_n\), 对应于\(X_K\)上的一个凝聚层集合\(\mathscr {U}(\mathscr {F}_1,\ldots ,\mathscr {F}_n)\), 并具有与层的同构及基底变换下的逆像相容这一显然的性质, 则给它配以一个仍记作\(\mathscr {U}\)的运算, 把\(n\)个凝聚层的类的集合所组成的序列\(E_1,\ldots ,E_n\)对应于一个凝聚层的类的集合\(\mathscr {U}(E_1,\ldots ,E_n)\).
\par{}本节的目的是定义某些称为\emph{有界}的层的类的集合, 并证明最常见的运算\(\mathscr {U}\)施于有界集合时, 所得集合仍然有界. 设\(X\)是在\(S\)上有限型的预概形, 且\(S\)为\emph{Noether}. 对每个\(s\in  S\), 纤维\(X_s\)是在\(k(s)\)上有限型的预概形; 我们考虑上述意义下\(X_s\)上的凝聚层的类. 这样,"\(X/S\)的一个纤维上的凝聚层的类"及类似说法便有了意义. 同样, 在各个纤维上分别进行上述步骤, 仍可考虑\(E\otimes  E'\)等运算: 它们把\(X/S\)的纤维上的凝聚层的类的集合族对应于\(X/S\)的纤维上的凝聚层的类的集合.
\begin{definition}{1.1}\setnumber{1.1}\label{fga3.iv-1-definition-1.1}\skipsinglepar
\par{}设\(E\)是\(X/S\)的纤维上的凝聚层的类的一个集合. 若存在\(S\)上的有限型预概形\(S'\)和\(X'=X\times _S S'\)上的凝聚层\(\mathscr {F}'\), 使\(E\)包含于\(\mathscr {F}'\)所定义的\(X/S\)的纤维上的层的类的集合中, 就称\(E\)是\emph{有界}的.
\end{definition}
\par{}按定义, 这个构造给\(s\in  S\)配以各层\(\mathscr {F}'\otimes _{S'}k(s')\)的类, 其中\(s'\)遍历\(S'\)中位于\(s\)之上的点 (于是\(k(s')\)是\(k(s)\)的扩张, 而\(X\otimes  (X_s)_{k(s')}\)可与\(X'\)在\(s'\)处的纤维\(X'\otimes _{S'}k(s')=X'_{s'}\)认同). 也可以说, 有界集合就是那些包含于某个凝聚层的\emph{代数族}中的集合, 此代数族由\(S\)上有限型的某个\(S'\)参数化.
\par{}有限多个有界集合的并有界 (取定义各个界定代数族的参数预概形\(S_i\)的和即可). 基底变换\(T\to  S\)把关于\(X/S\)有界的族变为关于\(X_T/T\)有界的族; 若\(T\to  S\)是满的, 则逆命题也成立 (更一般地, 只需它的像包含\(X/S\)的给定族\(E\)中出现的各个\(s\)). 因此在理论上, 只需在\(S\)为整数环\(\mathbb {Z}\)上有限型代数的谱的情形确定有界族.
\par{}若\(E\)和\(E'\)是关于\(X/S\)的层的类的有界族, 则\(E\otimes  E'\)也有界: 事实上, 若\(E\) (分别为\(E'\)) 被由\(T\to  S\)及\(X_T\)上的\(\mathscr {F}\) (分别为\(T'\to  S\)及\(X_{T'}\)上的\(\mathscr {F}'\)) 定义的代数族所界定, 则\(E\otimes  E'\)被下述代数族界定: 取\(T''=T\times _S T'\to  S\), 并在\(X_{T''}\)上取\(\mathscr {F}\)和\(\mathscr {F}'\)在\(X_{T''}\)上的逆像的张量积为层\(\mathscr {F}''\). 此论证成立, 正是因为函子\(\mathscr {F}\otimes \mathscr {F}'\)对\(\mathscr {F}\)和\(\mathscr {F}'\)均右正合, 因而与基底扩张交换, 特别与取纤维交换. 它不能原样用于\(\mathscr {T}\kern  -2.75pt\mathit {or}_i(E,E')\)、\(\mathscr {H}\kern  -2.5pt\mathit {om}(E,E')\)、\(\mathscr {E}\kern  -1.5pt\mathit {xt}^i(E,E')\)等局部运算. 不过, 按照\(E\otimes  E'\)的处理方式, 再利用下面这类结果 (均见[\hyperref[Gro1960b]{3}, IV 6.11]), 就可证明这些运算也把有界集合变为有界集合: 有界族\(E\)总能被某个\(X_T\)上的凝聚层\(\mathscr {F}\)所定义的代数族界定, 其中\(T\)在\(S\)上有限型, 且该层关于\(T\)\emph{平坦}. (将原来的参数空间"分成若干块".) 对适当的层要求这种平坦性, 确能保证\(\mathscr {T}\kern  -2.75pt\mathit {or}_i(\mathscr {F},\mathscr {F}')\)等运算与任意基底变换交换. 同一方法也适用于全局性质的运算: 固有态射下凝聚层的正像和导出正像、关于固有态射的全局\(\operatorname {Ext}\) (参见[\hyperref[GD1960]{5}, III §6]) 等; 所有这些运算都把层的有界族变为层的有界族 (注意: 这些运算可能改变承载所讨论各层的预概形).
\par{}下面两个断言可以用\emph{本质上}相同的平坦性方法证明; 关于有限型态射的纤维上的准素分解, 特别参见[\hyperref[GD1960]{5}, IV].
\begin{proposition}{1.2}\setnumber{1.2}\label{fga3.iv-1-proposition-1.2}\skipsinglepar
\par{}设\(E\)和\(E'\)是\(X/S\)的纤维上的层的类的有界集合, 并假设\(X\)在\(S\)上固有. 则
\begin{enumerate}
      
 \item[i.]同态\(\mathscr {F}\to \mathscr {F}'\)的核、余核和像所成的族有界 (其中\(\mathscr {F}\)的类属于\(E\), \(\mathscr {F}'\)的类属于\(E'\));

      
 \item[ii.]\(\mathscr {F}\)由\(\mathscr {F}'\)的扩张\(\mathscr {F}''\)所成的族有界 (其中\(\mathscr {F}\)的类属于\(E\), \(\mathscr {F}'\)的类属于\(E'\)).

    
\end{enumerate}
\begin{proof}
\par{}必要时作适当的基底变换, 可以假设\(E\)和\(E'\)分别由某个\(X_T/T\)上的凝聚层\(\mathscr {G}\)和\(\mathscr {G}'\)定义, 其中\(T\)在\(S\)上有限型. 还可以假设满足某些平坦性条件, 从而层\(f_T(\mathscr {H}\kern  -2.5pt\mathit {om}_{\mathcal {O}_{X_T}}(\mathscr {G},\mathscr {G}'))\)及\(\mathscr {E}\kern  -1.5pt\mathit {xt}_{f_T}^1(\mathscr {G},\mathscr {G}')\)的构造与任意态射\(T'\to  T\)所给出的基底变换交换. 又可假设上述凝聚层在\(T\)上局部自由. 设\(T_0\)、\(T_1\)为\(T\)上的向量丛, 它们的截面芽层分别是上述两层. 此时可以典范地定义\(X_{T_0}\)上的凝聚层同态\(\mathscr {G}_{T_0}\to \mathscr {G}'_{T_0}\), 以及\(X_{T_1}\)上的凝聚层扩张\[         0 \to  \mathscr {G}_{T_1} \to  \mathscr {G}'' \to  \mathscr {G}'_{T_1} \to  0,\] 后者具有显然的泛性质. 这个扩张中的层定义了一个代数族, 它界定(ii)中的族. 上述同态的核、余核和像也有同样性质; 考虑它们就证明(i) (这里须假设余核关于\(T_0\)平坦, 而这仍可归结到把\(T_0\)分块的办法...).
\end{proof}
\end{proposition}
\begin{proposition}{1.3}\setnumber{1.3}\label{fga3.iv-1-proposition-1.3}\skipsinglepar
\par{}设\(E\)是\(X/S\)的纤维上的层的类的有界族. 则各个\((\operatorname {supp}\mathscr {F})_\mathrm {red}\)的结构层的类组成一个有界族, 其中\(\mathscr {F}\)是某个\(X_K\)上的凝聚层, \(K\)代数闭, 且其类属于\(E\).
\end{proposition}
\par{}这里\((\operatorname {supp}\mathscr {F})_\mathrm {red}\)表示\(\mathscr {F}\)的支集, 赋以诱导的约化结构; 也就是说, 其结构层是\(\mathcal {O}_{X_K}\)模掉定义\(\operatorname {supp}\mathscr {F}\)的最大理想层所得的商. 对由准素分解理论从\(\mathscr {F}\)典范地导出的层, 可以证明与\hyperref[fga3.iv-1-proposition-1.3]{命题~\ref*{fga3.iv-1-proposition-1.3}}相类似的结果; 例如\(\mathscr {F}/\mathscr {F}_i\), 其中\(\mathscr {F}_i\)是\(\mathscr {F}\)中对应于\(\mathscr {F}\)的支集的各分支、并且关于此性质极小的准素子层; 又例如\(\mathcal {O}_{X_K}/\mathfrak {p}\), 其中\(\mathfrak {p}\)是\(\mathscr {F}\)的一个相伴素理想层; 或\(\mathcal {O}_{X_K}/\mathfrak {q}\), 其中\(\mathfrak {q}\)是对应于\(\mathscr {F}\)的支集的某个分支的准素理想层 (所取的基域为代数闭域).
\chapter{有界族与Hilbert多项式}\thispagestyle{fancy}\label{fga3.iv-2}
\par{}以下假设\(X\)在\(S\)上射影, 并带有一个非常丰沛层, 记作\(\mathcal {O}_X(1)\). 对\(S\)的一点\(s\)的剩余域\(k(s)\)的每个扩张\(K\), 考虑\(X_K\)上相应的层\(\mathcal {O}_{X_K}(1)\), 它仍然非常丰沛.
\par{}对\(X_K\)上每个凝聚层\(\mathscr {F}\), 配以函数\[     P_\mathscr {F}(n)     = \text {层 }\mathscr {F}(n)\text {在 }X_K\text { 上的Euler–Poincaré特征数}.\] 它是关于整数\(n\)的多项式, 称为\(\mathscr {F}\)的\emph{Hilbert多项式}. 对充分大的\(n\), \(P(n)\)恰为\(\operatorname {H}^0(X_K,\mathscr {F}(n))\)在\(K\)上的维数, 因为对\(i>0\)和充分大的\(n\), 各\(\operatorname {H}^i(X_K,\mathscr {F}(n))\)为零.
\par{}若\(\mathscr {F}\)是\(X\)上关于\(S\)\emph{平坦}的凝聚层, 则在纤维\(X_s\)上诱导的层\(\mathscr {F}_s\)的Hilbert多项式 (对于\(S\)的同一个连通分支) 均相同 [\hyperref[GD1960]{5}, III, §7]. 因而对\(X\)上的每个凝聚层\(\mathscr {F}\), 各层\(\mathscr {F}_s\) (\(s\in  S\)) 的Hilbert多项式的集合为有限集, 这里不需要平坦性假设.
\par{}还须回顾, \(X\)上的任意凝聚层\(\mathscr {F}\)都同构于某个形如\(\mathcal {O}_X(-n)^N\)的层的商层, 其中\(n,N\)充分大. 因而纤维上诱导的层\(\mathscr {F}_s\)也都是纤维上相应的层\(\mathcal {O}(-n)\)的商.
\par{}从这两点得到下述定理中条件的必要性:
\begin{theorem}{2.1}\setnumber{2.1}\label{fga3.iv-2-theorem-2.1}\skipsinglepar
\par{}设\(X\)在\(S\)上射影, \(S\)为Noether, 且\(\mathcal {O}_X(1)\)是\(X\)上关于\(S\)非常丰沛的层. 设\(E\)是\(X/S\)的纤维上的层的类的一个集合. \(E\)有界的充要条件是满足下列条件: 
\begin{enumerate}
      
 \item[a.]\(X\)上存在一个凝聚层\(\mathscr {L}\) (可以假设形如\(\mathcal {O}_X(-n)^N\)), 使得\(E\)包含于形如\(\mathscr {L}_K\)的层的凝聚商层的类所组成的族;

      
 \item[b.]类属于\(E\)的各层\(\mathscr {F}\)的Hilbert多项式\(P_\mathscr {F}\), 都属于同一个\emph{有限}多项式集合.

    
\end{enumerate}
\end{theorem}
\par{}尚须证明条件的充分性, 这将是一个更精确结果的特例. 对域\(K\)上有限型预概形上的每个凝聚模\(\mathscr {F}\), 以及每个整数\(r\), 以\(N_r\)表示\(\mathscr {F}\)的一个子模: 在任一开子集上的截面, 是\(\mathscr {F}\)在该开子集上支集维数\(<r\)的截面. 于是对\(r>\dim \operatorname {supp}\mathscr {F}\)有\(N_r=\mathscr {F}\), 对\(r\leqslant 0\)有\(N_r=0\), 从而得到\(\mathscr {F}\)的一个有限递增滤过, 其因子\(N_r/N_{r+1}\)的相伴素闭链, 恰为\(\mathscr {F}\)的维数为\(r\)的相伴素闭链. 置\[     \mathscr {F}_{(r)}     = \mathscr {F}/N_r,\] 则\(\mathscr {F}_{(r)}\)的相伴素闭链恰为\(\mathscr {F}\)的维数\(\geqslant  r\)的相伴素闭链; 特别地, \(\mathscr {F}_{(r)}\)等于\(\mathscr {F}\)当且仅当\(\mathscr {F}\)的相伴素闭链的维数均\(\geqslant  r\). 此时有:
\begin{theorem}{2.2}\setnumber{2.2}\label{fga3.iv-2-theorem-2.2}\skipsinglepar
\par{}在\hyperref[fga3.iv-2-theorem-2.1]{定理~\ref*{fga3.iv-2-theorem-2.1}}的条件下, 设\(s\)是整数, 并假设\(E\)满足条件(a)以及(b)的下述弱化形式: 
\begin{enumerate}
      
 \item[bs.]类属于\(E\)的各层\(\mathscr {F}\)的Poincaré多项式\(P_\mathscr {F}\), 其\emph{次数\(\leqslant  (s-1)\)的各项}的系数有界.

    
\end{enumerate} 此时, 各层\(\mathscr {F}_{(s)}\) (其中\(\mathscr {F}\)的类属于\(E\)) 组成一个有界族. 此外, \(P_\mathscr {F}\)的\((s-2)\)次项系数均有下界.
\end{theorem}
\par{}于是:
\begin{corollary}{2.3}\setnumber{2.3}\label{fga3.iv-2-corollary-2.3}\skipsinglepar
\par{}假设类属于\(E\)的各层\(\mathscr {F}\), 其所有相伴素闭链的维数\(d\)均满足\(s\leqslant  d\leqslant  r\). 那么在\hyperref[fga3.iv-2-theorem-2.1]{定理~\ref*{fga3.iv-2-theorem-2.1}}的条件(b)中, 只需考虑\(P_\mathscr {F}\)的从\((s-1)\)次到\(r\)次 (两端包括在内) 的系数.
\end{corollary}
\par{}本节的其余部分用来概述\hyperref[fga3.iv-2-theorem-2.2]{定理~\ref*{fga3.iv-2-theorem-2.2}}的证明. 关键是下面两个引理, 其中第一个是熟知的 (它概括了Chow坐标所包含的有效数学内容):
\begin{lemma}[Chow]{2.4}\setnumber{2.4}\label{fga3.iv-2-lemma-2.4}\skipsinglepar
\par{}考虑纤维\(X_K\)的子概形\(Y\)的结构层, 其中\(K\)是\(S\)的某个剩余域的代数闭扩张, \(Y\)约化且所有分支具有相同维数\(r\) (并把\(\mathcal {O}_Y\)视为\(\mathcal {O}_X\)的商层). 若\(Y\)的次数有界, 则这些\(Y\)组成一个有界族.
\end{lemma}
\par{}这里\(Y\)的次数\(a\)最方便地定义为\(P_{\mathcal {O}_Y}=an^r/r!+\ldots \)中最高次项的系数.
\begin{lemma}{2.5}\setnumber{2.5}\label{fga3.iv-2-lemma-2.5}\skipsinglepar
\par{}设\(\mathscr {L}\)是\(X\)上的凝聚层, \(E\)是层\(\mathscr {L}_K\)的商层\(\mathscr {F}\)的类所组成的一个集合 (其中\(K\)是\(S\)的一个剩余域的扩张). 假设\(X\)在\(S\)上的各纤维的维数\(\leqslant  r\), 并置\[       P_\mathscr {F}(n)       = a_\mathscr {F}n^r/r! + b_\mathscr {F}n^{r-1}/(r-1)! + \text {次数}<r-1\text { 的项}.\] 则系数\(a_\mathscr {F}\)有 (上) 界, 而\(b_\mathscr {F}\)有下界. 若\(b_\mathscr {F}\)有界, 则族\(\mathscr {F}_{(r)}\)有界.
\begin{proof}
\par{}将\(S\)替换为覆盖\(S\)的\(S\)的若干子概形的并, 可以假设存在\emph{有限}态射\(f\colon  X\to \mathbb {P}_S^r\), 使\(\mathcal {O}_X(1)\)同构于\(\mathcal {O}_{\mathbb {P}_S^r}(1)\)的逆像; 因而对于\(X\)上的每个凝聚层\(\mathscr {F}\), 均有\(P_\mathscr {F}=P_{f_*(\mathscr {F})}\). 利用上一节的方法, 也容易证明\(X\)上的层\(\mathscr {F}\)的集合有界, 当且仅当各\(f_*(\mathscr {F})\)的集合有界. 最后, 有\[         f_*(\mathscr {F})_{(r)}         = f_*(\mathscr {F}_{(r)}).\] 因此可归结到\(X=\mathbb {P}_S^r\)的情形. 还可对适当的\(k\)和\(s\)假设\(\mathscr {L}=\mathcal {O}_{\mathbb {P}_S^r}(k)^s\). 系数\(a_\mathscr {F}\)满足\[         0\leqslant  a_\mathscr {F} \leqslant  s,\] 故有界. 据此, 说\(P_\mathscr {F}(n)\)的\(n^{r-1}\)次项系数有下界 (分别为有界), 等价于对\(P_\mathscr {F}(n-k)=P_{\mathscr {F}(-k)}(n)\)作同样断言. 这样就归结到\[         \mathscr {L}         = \mathcal {O}_{\mathbb {P}_S^r}^s.       \]的情形.
\par{}考虑正合列\[         0 \to  N_r \to  \mathscr {F} \to  \mathscr {F}_{(r)} \to  0,\] 由此得\[         P_\mathscr {F}         = P_{\mathscr {F}_{(r)}} + P_{N_r}.\] 又因\(P_{N_r}\)的\(n^{r-1}\)次项系数非负 (因为\(\dim \operatorname {supp} N_r\leqslant  r-1\)), 所以\[         b_{\mathscr {F}_{(r)}} \leqslant  b_\mathscr {F}.\] 因此在证明引理时, 可以用\(\mathscr {F}_{(r)}\)替代\(\mathscr {F}\), 即假设所讨论的\(\mathscr {L}\)的商\(\mathscr {F}\)无挠.
\par{}由于\(\mathbb {P}_K^r\)正规, 所以\(\mathscr {F}\)在某个开集\(U=\mathbb {P}_K^r\setminus  Y\)上是秩为\(a=a_\mathscr {F}\)的局部自由层, 其中\(Y\)的余维\(\geqslant 2\). 因而\(\bigwedge ^a\mathscr {F}\)是\(\mathbb {P}_K^r\)上的一个层, 其在\(U\)上的限制可逆, 从而同构于\(\mathbb {P}_K^r\)上的一个在相差同构意义下确定的可逆层的限制 (因为\(\mathbb {P}_K^r\)正则, 且\(Y\)的余维\(\geqslant 2\)). 后一个层形如\(\mathcal {O}_{\mathbb {P}_K^r}(d)\), 其中\(d\)是完全确定的整数. 由于\(\bigwedge ^a\mathscr {F}\)是\(\bigwedge ^a\mathcal {O}_{\mathbb {P}_K^r}^n\simeq \mathcal {O}_{\mathbb {P}_K^r}^N\)的一个商, 其中\(N=\binom {n}{a}\), 所以它有\(N\)个典范截面. 这些截面给出\(\mathcal {O}_{\mathbb {P}_K^r}(d)\)在\(U\)上的截面, 后者是\(\mathcal {O}_{\mathbb {P}_K^r}(d)\)的截面\(s_i\) (\(1\leqslant  i\leqslant  N\)) 的限制 (因为\(\mathbb {P}_K^r\)正规, 且\(Y\)的余维\(\geqslant 2\)). 这些\(s_i\)在\(U\)的各点生成\(\mathcal {O}_{\mathbb {P}_K^r}(d)\), 因此不全为零, 从而\(d\geqslant 0\). 一个容易的计算还给出\[         b_\mathscr {F}         = a_\mathscr {F}(r+1)/2 + d.\] 特别有\(b_\mathscr {F}\geqslant 0\), 所以\(b_\mathscr {F}\)有下界. 它有界当且仅当\(d\)有界; 我们将证明此时\(\mathscr {F}\)仍属于一个有界族. 可以固定\(a_\mathscr {F}\)和\(b_\mathscr {F}\), 记为\(a\)和\(b\) (从而\(d\)也固定). 给定\(\mathcal {O}_{\mathbb {P}_K^r}(d)\)的\(N\)个截面\(s_i\), 即给定一个同态\(s\colon \bigwedge ^a\mathscr {L}_K\to \mathcal {O}_{\mathbb {P}_K^r}(d)\), 就能把\(\mathscr {F}\)恢复为相应复合同态\[         \mathscr {L}_K         \to  \mathscr {H}\kern  -2.5pt\mathit {om}\left ( \bigwedge ^{a-1}\mathscr {L}_{K'}, \bigwedge ^a\mathscr {L}_K \right )         \to  \mathscr {H}\kern  -2.5pt\mathit {om}\left ( \bigwedge ^{a-1}\mathscr {L}_{K'}, \mathcal {O}_{\mathbb {P}_K^r}(d) \right )       \]的余像; 其中第一个箭头是外积给出的典范同态, 第二个箭头来自\(s\). 再由\hyperref[fga3.iv-1-proposition-1.2]{命题~\ref*{fga3.iv-1-proposition-1.2}}的(i)即得结论.
\end{proof}
\end{lemma}
\par{}合用上述两个引理, 可以证明:
\begin{lemma}{2.6}\setnumber{2.6}\label{fga3.iv-2-lemma-2.6}\skipsinglepar
\par{}在\hyperref[fga3.iv-2-theorem-2.1]{定理~\ref*{fga3.iv-2-theorem-2.1}}的初始条件下, 假设对所有\(\mathscr {F}\)都有\[       P_\mathscr {F}(n)       = a_\mathscr {F} n^r/r! + b_\mathscr {F} n^{r-1}/(r-1)! + \text {次数}<r-1\text { 的项},\] 且系数\(a_\mathscr {F}\)有界. 则系数\(b_\mathscr {F}\)有下界. 此外, 若各\(b_\mathscr {F}\)有界, 则各\(\mathscr {F}_{(r)}\)有界.
\begin{proof}
\par{}可以假设各层\(\mathscr {F}\)的基域\(K\)代数闭. 给每个\(\operatorname {supp}\mathscr {F}_{(r)}\) (即\(r\)维分支的并) 赋以诱导的约化结构. 则各\(\operatorname {supp}\mathscr {F}_{(r)}\)的次数以\(a\)为上界, 所以由\hyperref[fga3.iv-2-lemma-2.4]{引理~\ref*{fga3.iv-2-lemma-2.4}}, 各\(\operatorname {supp}\mathscr {F}_{(r)}\)组成有界集合. 又对于\(\operatorname {supp}\mathscr {F}_{(r)}\)的每个分支, \(\mathscr {F}_{(r)}\)在该分支处的长度\(\leqslant  a\); 因而若\(\mathscr {I}_\mathscr {F}\)是定义\(\operatorname {supp}\mathscr {F}_{(r)}\)的理想, 则可将\(\mathscr {F}_{(r)}\)视为\(X\)中由\(\mathscr {I}_\mathscr {F}^a\)定义的子概形\(Y_\mathscr {F}\)上的模. 如上一引理, 还可归结到\(\mathscr {F}=\mathscr {F}_{(r)}\)的情形, 从而\(\mathscr {F}\)来自\(Y_\mathscr {F}\)上的一个模. 各\(Y_\mathscr {F}\)对应于\(\mathcal {O}_{X_K}\)的商模的一个有界族, 因而来自某个概形\(X\times _S T\)的闭子概形\(Y\). 于是对\(Y/T\)及\(\mathscr {L}\otimes _X Y\)应用\hyperref[fga3.iv-2-lemma-2.5]{引理~\ref*{fga3.iv-2-lemma-2.5}}, 即得结论.
\end{proof}
\end{lemma}
\par{}现在对\(\dim \operatorname {supp}\mathscr {F}\)的上界\(r\)作归纳, 证明\hyperref[fga3.iv-2-theorem-2.2]{定理~\ref*{fga3.iv-2-theorem-2.2}}. 当\(r<0\)时断言显然成立; 假设\(r\geqslant 0\), 且对\(r'<r\)已证得断言. 由\hyperref[fga3.iv-2-lemma-2.6]{引理~\ref*{fga3.iv-2-lemma-2.6}}, 各\(\mathscr {F}_{(r)}\)组成一个有界族; 再由\hyperref[fga3.iv-1-proposition-1.2]{命题~\ref*{fga3.iv-1-proposition-1.2}}的(i), 各同态\(\mathscr {L}_K\to \mathscr {F}_{(r)}\)的核也组成有界族. 因而\(X\)上存在一个凝聚模\(\mathscr {L}'\), 使这些核, 进而也使各\(N_r(\mathscr {F})=\operatorname {Ker}(\mathscr {F}\to \mathscr {F}_{(r)})\), 都是模\(\mathscr {L}'_K\)的商. 因为各\(\mathscr {F}_{(r)}\)有界, 所以各\(P_{\mathscr {F}_{(r)}}\)有界, 公式\[     P_\mathscr {F}     = P_{\mathscr {F}_{(r)}} + P_{N_r}   \]便说明各\(P_{N_r}\)满足与各\(P_\mathscr {F}\)相同的条件(bs). 因而各\(N_r\)满足条件(a)和(bs), 由归纳假设, 各\((N_r)_{(s)}\)有界. 而\(\mathscr {F}_{(s)}\)是\(\mathscr {F}_{(r)}\)由\((N_r)_{(s)}\)的扩张, 所以由\hyperref[fga3.iv-1-proposition-1.2]{命题~\ref*{fga3.iv-1-proposition-1.2}}的(ii), 各\(\mathscr {F}_{(s)}\)有界. 对于\hyperref[fga3.iv-2-theorem-2.2]{定理~\ref*{fga3.iv-2-theorem-2.2}}的最后一个断言, 注意由\hyperref[fga3.iv-1-proposition-1.2]{命题~\ref*{fga3.iv-1-proposition-1.2}}的(i), 各映射\(\mathscr {F}\to \mathscr {F}_{(s)}\)的核\(N_s\)有界, 且\(P_{N_s}\)的\(n^{s-1}\)次项系数有界; 此时\hyperref[fga3.iv-2-lemma-2.6]{引理~\ref*{fga3.iv-2-lemma-2.6}}说明下一项的系数有下界. \hyperref[fga3.iv-2-theorem-2.2]{定理~\ref*{fga3.iv-2-theorem-2.2}}的证明至此完成.
\chapter{Hilbert概形: 定义与存在性定理}\thispagestyle{fancy}\label{fga3.iv-3}
\par{}设\(X\)是预概形\(S\)之上的预概形, \(\mathscr {F}\)是\(X\)上的拟凝聚模. 用\[     \operatorname {Quot}({\mathscr {F}/X/S})   \]表示\(\mathscr {F}\)的各个在\(S\)上平坦的拟凝聚商模所组成的集合. 现设\(S'\to  S\)为基底变换态射, 置\(X'=X\times _S S'\)和\(\mathscr {F}'=\mathscr {F}\otimes _{\mathcal {O}_S}\mathcal {O}_{S'}\); 则\(X'\)是\(S'\)上的预概形, 带有拟凝聚模\(\mathscr {F}'\), 因而可以考虑\(\operatorname {Quot}({\mathscr {F}'/X'/S'})\). 置\[     \mathscr {Q}\kern  -.75pt\mathit {uot}_{{\mathscr {F}/X/S}} (S')     = \operatorname {Quot}({\mathscr {F}'/X'/S'})   \] (其中仍如上置\(X'=X\times _S S'\)).
\par{}若\(S''\to  S'\)是一个\(S\)-态射, 则\(X''=X\times _S S''\)同构于\(X'\times _{S'}S''\), \(\mathscr {F}''\)同构于\(\mathscr {F}'\otimes _{\mathcal {O}_{S'}}\mathcal {O}_{S''}\). 因为逆像函子\(\mathscr {G}'\mapsto \mathscr {G}'\otimes _{\mathcal {O}_{S'}}\mathcal {O}_{S''}\)右正合, 取值于\(X''\)上的拟凝聚模范畴, 并把\(S'\)-平坦模变成\(S''\)-平坦模, 所以得到自然映射\[     \operatorname {Quot}({\mathscr {F}'/X'/S'}) \to  \operatorname {Quot}(\mathscr {F}''/X''/S''),\] 从而\(\mathscr {Q}\kern  -.75pt\mathit {uot}_{{\mathscr {F}/X/S}} (S')\)是一个取值于集合范畴的\emph{关于\(S'\)的反变函子} (其中\(S'\)是\(S\)上的预概形). 以下假设\(X\)在\emph{Noether}的\(S\)上\emph{射影}, 且\(\mathscr {F}\)为\emph{凝聚}的; 为简单起见, 只考虑在\(S\)上的\emph{局部Noether}的\(S'\).
\begin{theorem}{3.1}\setnumber{3.1}\label{fga3.iv-3-theorem-3.1}\skipsinglepar
\par{}在这些条件下, 局部Noether的\(S\)-预概形范畴上的反变函子\(\mathscr {Q}\kern  -.75pt\mathit {uot}_{{\mathscr {F}/X/S}}\), 可由一个\(S\)-预概形\(\underline {\operatorname {Quot}}_{{\mathscr {F}/X/S}}\)\emph{表示}, 这个预概形是一个射影\(S\)-概形序列的和 (因而\(\underline {\operatorname {Quot}}_{{\mathscr {F}/X/S}}\)在\(S\)上局部有限型).
\end{theorem}
\par{}我们将按以下方式得到这样的分解. 设\(\mathcal {O}_X(1)\)是\(X\)上关于\(S\)非常丰沛的可逆层. 对每个有理系数多项式\(P(n)\), 以\(\operatorname {Quot}^P({\mathscr {F}/X/S})\)表示\(\operatorname {Quot}({\mathscr {F}/X/S})\)的一个子集, 其元素为\(\mathscr {F}\)的凝聚商\(\mathscr {G}\), 它在\(S\)上平坦, 且在每个\(s\in  S\)处的Hilbert多项式均为\(P\). 再置\[     \mathscr {Q}\kern  -.75pt\mathit {uot}_{\mathscr {F}/X/S}^P(S')     = \operatorname {Quot}^P({\mathscr {F}'/X'/S'}),\] 就得到\(\mathscr {Q}\kern  -.75pt\mathit {uot}_{\mathscr {F}/X/S}\)的一个子函子. §2所回顾的Hilbert多项式的不变性质给出: \emph{\(\mathscr {Q}\kern  -.75pt\mathit {uot}_{\mathscr {F}/X/S}\)可表示的充要条件是各\(\mathscr {Q}\kern  -.75pt\mathit {uot}_{\mathscr {F}/X/S}^P\)可表示; 此时表示它的\(S\)-预概形\(\underline {\operatorname {Quot}}_{{\mathscr {F}/X/S}}\)同构于表示各函子\(\mathscr {Q}\kern  -.75pt\mathit {uot}_{\mathscr {F}/X/S}^P\)的预概形\(\underline {\operatorname {Quot}}_{{\mathscr {F}/X/S}}^P\)之和.} 据此, \hyperref[fga3.iv-3-theorem-3.1]{定理~\ref*{fga3.iv-3-theorem-3.1}}将由下述定理推出:
\begin{theorem}{3.2}\setnumber{3.2}\label{fga3.iv-3-theorem-3.2}\skipsinglepar
\par{}在上述记号下, 函子\(\mathscr {Q}\kern  -.75pt\mathit {uot}_{\mathscr {F}/X/S}^P\)可由一个射影\(S\)-预概形\(\operatorname {Quot}_{\mathscr {F}/X/S}^P\)表示.
\end{theorem}
\par{}本节其余部分证明\hyperref[fga3.iv-3-theorem-3.2]{定理~\ref*{fga3.iv-3-theorem-3.2}}.
\par{}设\(\nu \)是整数. 对每个\(S\)上的\(S'\), 以\(A_\nu (S')\)表示\(\mathscr {F}'=\mathscr {F}\otimes _{\mathcal {O}_S}\mathcal {O}_{S'}\)的各个凝聚商\(\mathscr {G}=\mathscr {F}'/\mathscr {H}\)所组成的集合, 其中要求这些商在\(S'\)上平坦, 且满足下列条件: 
\begin{enumerate}
    
 \item[a.]当\(i>0\)且\(n\geqslant \nu \)时, \(\operatorname {R}^i f'_*(\mathscr {G} (n))=0\);

    
 \item[b.]当\(i>0\)且\(n\geqslant \nu \)时, \(\operatorname {R}^i f'_*(\mathscr {H} (n))=0\);

    
 \item[c.]当\(k\geqslant 0\)时, \(f'_*(\mathscr {H} (\nu +k))=\mathcal {S}'_k f'_*(\mathscr {H} (\nu ))\).

  
\end{enumerate} 对最后一个条件, 假设\(X\)写成\(S\)上一个由\(\mathcal {S}_1\)生成的拟凝聚正分次代数\(\mathcal {S}_\bullet \)的齐次素谱, 并置\(\mathcal {S}'=\mathcal {S}\otimes _{\mathcal {O}_S}\mathcal {O}_{S'}\), 于是\(X'\)是\(\mathcal {S}'\)的齐次素谱. 为证明\hyperref[fga3.iv-3-theorem-3.2]{定理~\ref*{fga3.iv-3-theorem-3.2}}, 很容易归结到\(X=\mathbb {P}_S^r\)的情形 (因为\(S\)是开子集\(U\)的并, 在这些开子集上\(X|U\)是\(\mathbb {P}_U^r\)的闭子概形, 且\(\mathcal {O}_X(1)\)由\(\mathcal {O}_{\mathbb {P}_U^r} (1)\)诱导), 并归结到\(\mathscr {F}\)形如\((\mathcal {O}_{\mathbb {P}_S^r})^N\), 因而在\(S\)上平坦的情形. 此时, 上述各层\(\mathscr {H}\)也在\(S'\)上平坦. 于是由Künneth关系 [\hyperref[GD1960]{5}, III, §7] 以及(b), 条件(a)和(b)在基底变换下稳定, 并且蕴涵: 当\(n\geqslant \nu \)时, \(f'_*(\mathscr {G} (n))\)和\(F'_*(\mathscr {H} (n))\)的构造与基底扩张交换 (同上), 因而(c)也在基底扩张下稳定. 换言之, \emph{\(A_\nu (S')\)是关于\(S'\)的反变函子, 并且在准确的意义下是\(A(S')=\mathscr {Q}\kern  -.75pt\mathit {uot}_{{\mathscr {F}/X/S}}^P(S')\)的子函子}. 令\(\nu \)变动, 就得到\(A(S')\)的一列递增子集\(A_\nu (S')\); 由Serre的一个熟知定理 [\hyperref[GD1960]{5}, III, §2], 其并为\(A(S')\). 现注意, 若\(\mathscr {G}\)是\(\mathscr {F}'\)的凝聚商, 在\(S\)上平坦, 而\(S\)的一个元素\(s\)使得基底变换\(\operatorname {Spec}(k(s))\to  S\)所得的\(\mathscr {F}_s\)的商\(\mathscr {G}_s\)满足(a)、(b)、(c), 即属于\(A_\nu \operatorname {Spec}(k(s))\), 则\(s\)有一个开邻域\(U\), 使\(\mathscr {G}|(f')^{-1} (U)\)也满足同样条件, 即这个商属于\(A_\nu (U)\); 对(a)和(b), 这事实上由"全纯函数定理" [\hyperref[GD1960]{5}, III, §7] 得到; 对(c), 则由Nakayama引理以及对充分大的\(n\)和\(k\geqslant 0\)总有\(f'_*(\mathscr {H} (n+k))=S'_kf'_*(\mathscr {H} (n))\)这一事实 [\hyperref[GD1960]{5}, III, §2] 得到. 由这些说明可得以下结论 (比较[\hyperref[Gro1960a]{4}, IV]): \emph{函子\(A\)可表示的充要条件是各函子\(A_\nu \)可表示; 此时表示\(A\)的\(S\)-预概形\(Q\), 是表示各\(A_\nu \)的开子集\(Q_\nu \)的递增并}
\par{}置\[     M_\bullet      = \sum _{n\geqslant 0} f_*(\mathscr {F} (n))     = \mathcal {S}_\bullet ^N,\] 因而有\[     M'_\bullet      = M_\bullet \otimes _{\mathcal {O}_S}\mathcal {O}_{S'}     = \sum _{n\geqslant 0} f'_*(\mathscr {F} (n))     = {\mathcal {S}'_\bullet }^N.\]
\par{}由(a)有
\begin{enumerate}
    
 \item[a'.]当\(n\geqslant \nu \)时, \(f'_*(\mathscr {G} (n))\)是秩为\(P(n)\)的局部自由层 [\hyperref[GD1960]{5}, III, §7]

  
\end{enumerate}, 而由(b)在\(i=1\)时有
\begin{enumerate}
    
 \item[a''.]\(f'_*(\mathscr {G} (n))\)是\(M'_n\)的商模.

  
\end{enumerate}. 又对\(n=\nu \)知道这个商模, 就由(c)知道\(M_n\)在\(n\geqslant \nu \)时的各个子模\(f'_*(\mathscr {H} (n))\), 从而知道\(\mathscr {H}\), 进而知道\(\mathscr {G}\). 这样得到一个\emph{单射}\[     A_\nu (S') \to  \mathscr {G}\kern  -1.25pt\mathit {rass}_{P(\nu )} (M'_\nu ),\] 从\(A_\nu (S')\)到\(M'\)的秩为\(P(\nu )\)的局部自由商模的集合; 因而得到一个\emph{函子性}的同态\[     i_\nu \colon  A_\nu (S') \to  \mathscr {G}\kern  -1.25pt\mathit {rass}_{P(\nu )} (M_\nu )(S').\] 右侧函子可由Grassmann概形\(\operatorname {Grass}_{P(\nu )} (M_\nu )\)表示 (比较[\hyperref[Gro1960a]{4}, V]), 而此概形在\(S\)上射影. 于是有:
\begin{lemma}{3.3}\setnumber{3.3}\label{fga3.iv-3-lemma-3.3}\skipsinglepar
\par{}\(A_\nu (S')\)是可表示函子, 且表示同态\(i_\nu \)的态射\(Q_\nu \to \operatorname {Grass}_{P(\nu )} (M_\nu )\)是浸入 (这蕴涵\(Q_\nu \)在\(S\)上拟射影).
\end{lemma}
\par{}这个断言等价于以下断言 (比较[\hyperref[Gro1960a]{4}, IV]): \emph{若给定\(M'_\nu \)的一个局部自由秩为\(P(\nu )\)的商模\(N\), 则\(S'\)有一个子预概形\(Z\), 使得对于每个\(S'\)上的局部Noether预概形\(T'\), \(N\)在\(T'\)上的逆像属于\(\Im  A_\nu (T')\)的充要条件是\(T'\to  S'\)经由子预概形\(Z\)分解.} 改换记号, 可以假设\(S'=S\), 即有\(M_\nu \)模掉子模\(R_\nu \)得到的商\(N_\nu \). 它来自\(A(s)\)的一个元素的充要条件是满足以下两个条件: 
\begin{enumerate}
    
 \item[i.]对\(k\geqslant 0\), \(M_{\nu +k}/\mathcal {S}_k R_\nu \)局部自由且秩为\(P(\nu +k)\).

    
 \item[ii.]由\(M_\bullet \)的分次子模\(R_\bullet =\sum _{k\geqslant 0}\mathscr {S}_kR_\nu \)定义的\(\mathscr {F}\)的子层\(\mathscr {H}\) (参见[\hyperref[GD1960]{5}, II, §3]) 及商层\(\mathscr {G}=\mathscr {F}/\mathscr {H}\)均满足上述条件(a)和(b).

  
\end{enumerate} 这些条件显然必要; 若它们满足, 则上述(ii)定义的层\(\mathscr {G}\)同构于分次\(\mathscr {S}_\bullet \)-模\(N_\bullet \) (即各\(M_{n+k}/\mathscr {S}_kR\)之和) 所对应的层, 所以在\(S\)上\emph{平坦} (因为它的各茎是\(N\)关于\(\mathcal {S}_\bullet \)的齐次素理想的局部化的直和因子), 而且由(i), 对应于Hilbert多项式\(P\). 再考虑(ii), 就看到(a')和(a'')满足; 因而当\(n\geqslant \nu \)时, 自然同态\(N_n\to  f_*(\mathscr {G} (n))\)是等秩局部自由模间的\emph{满}同态, 故为同构, 从而对所有\(k\geqslant 0\)有\(f_*(\mathscr {H} (\nu +k))=\mathcal {S}_kR_\nu \). 这就证明\(\mathscr {G}\in  A_\nu (S)\), 且\(R_\nu \)是由\(\mathscr {G}\)所定义的\(\operatorname {Grass}_{P(\nu )} (M_\nu )\)的元素.
\par{}上述判据(i)和(ii)同样适用于基底变换\(S'\to  S\)后的情形. 首先证明, 基底变换\(S'\to  S\)后条件(i)满足, 可以表示为\(S'\to  S\)经由\(S\)的某个子预概形\(Z\)分解. 一旦证明这一点, 用\(Z\)替代\(S\), 就归结到(i)已在\(S\)上满足的情形; 因为它在基底变换下稳定, 所以只须表述条件(ii). 但此时以\(U\)表示这样的\(s\in  S\)的集合: \(\mathscr {G} (n)\)与\(\mathscr {H} (n)\)在纤维\(X_s\)上诱导的层, 对\(n\geqslant \nu \)在所有\(>0\)维的上同调均为零. 我们已经证明\(U\)是开集, 而条件(ii)在基底变换\(S'\to  S\)后满足, 当且仅当\(S'\to  S\)经由\(U\)分解; 这就证明\hyperref[fga3.iv-3-lemma-3.3]{引理~\ref*{fga3.iv-3-lemma-3.3}}. 因此尚须证明以下引理:
\begin{lemma}{3.4}\setnumber{3.4}\label{fga3.iv-3-lemma-3.4}\skipsinglepar
\par{}设\(S\)是一个局部Noether预概形, 其上有一个由\(\mathcal {S}_1\)生成的拟凝聚正分次代数\(\mathcal {S}_\bullet \); 设\(M_\bullet \)是有限型拟凝聚分次\(\mathcal {S}\)-模, \(P\)是有理系数多项式, \(\nu \)是整数. 则\(S\)存在一个显然唯一的子预概形\(Z\), 具有如下性质: 对\(S\)上的每个预概形\(S'\), \(M_n\otimes _{\mathcal {O}_S}\mathcal {O}_{S'}\)在所有\(n\geqslant \nu \)时均局部自由且秩为\(P(n)\)的充要条件是\(S'\to  S\)经由\(Z\)分解.
\end{lemma}
\par{}显然可以假设\(S\)仿射, 因而为Noether. 于是:
\begin{lemma}{3.5}\setnumber{3.5}\label{fga3.iv-3-lemma-3.5}\skipsinglepar
\par{}对每个整数\(N\geqslant \nu \), 设\(U_N\)是\(S\)中由满足下列条件的\(s\in  S\)组成的开子集: 对所有\(\nu \leqslant  n\leqslant  N\), 均有\(\operatorname {rank}_{k(s)}M_{ns}\otimes _{\mathcal {O}_{S,s}} k(s)\leqslant  P(n)\). 则这列递减开子集\(U_N\)最终稳定.
\begin{proof}
\par{}已知 [\hyperref[Gro1960b]{3}, IV], \(S\)可以有限地分割成约化子概形\(S_i\), 使每个\(M\otimes _{\mathcal {O}_S}\mathcal {O}_{S_i}\)在\(S\)上\emph{平坦}. 因而可假设\(M\)平坦, 从而各\(M_n\)平坦. 最后, 显然可假设\(S\)连通. 但此时 [\hyperref[GD1960]{5}, III, §7], 存在一个整数\(n_0\)和一个多项式\(Q\), 使\[         \operatorname {rank}_{k(s)}M_{ns}\otimes _{\mathcal {O}_{S,s}} k(s)         = Q(n)         \qquad \text {当 } n\geqslant  n_0.\] 首先假设\(P<Q\), 从而对充分大的\(n\)有\(P(n)\neq  Q(n)\). 则显然对充分大的\(N\)有\(U_N=\varnothing \), 因而各\(U_N\)组成的序列当然稳定. 在相反情形, 对充分大的\(n\)有\(P(n)\geqslant  Q(n)\), 所以当\(N\geqslant  n_0\)时\(U_N=U_{n_0}\), 各\(U_N\)组成的序列同样稳定.
\end{proof}
\end{lemma}
\par{}特别地, 满足
\begin{eqn}{*}\setnumber{*}\label{fga3.iv-3-equation-star}\skipsinglepar\[       \operatorname {rank}_{k(s)}M_{ns}\otimes _{\mathcal {O}_{S,s}} k(s) \leqslant  P(n)       \qquad \text {对所有 } n\geqslant \nu      \tag{*}     \]
\end{eqn}的\(s\in  S\)所组成的集合\(U_\infty \)是开集, 因为它是各\(U_N\)之交. 因此在证明\hyperref[fga3.iv-3-lemma-3.4]{引理~\ref*{fga3.iv-3-lemma-3.4}}时, 可以用开子集\(U\)替代\(S\), 这就归结到\hyperref[fga3.iv-3-equation-star]{公式~\ref*{fga3.iv-3-equation-star}}中的不等式在所有\(s\in  S\)处成立的情形.
\begin{lemma}{3.6}\setnumber{3.6}\label{fga3.iv-3-lemma-3.6}\skipsinglepar
\par{}设\(M\)是局部Noether预概形\(S\)上的模, \(r\)是整数. 则\(S\)有一个显然唯一的子预概形\(Z\), 具有如下性质: 对每个\(S\)上的\(S'\), \(M\otimes _{\mathcal {O}_S}\mathcal {O}_{S'}\)局部自由且秩为\(r\)的充要条件是\(S'\to  S\)经由\(Z\)分解. 若对所有\(s\)都有\(\operatorname {rank}_{k(s)}M_s\otimes _{\mathcal {O}_{S,s}} k(s)\leqslant  r\), 则\(Z\)是\(S\)的闭子预概形 (假设\(M\)凝聚).
\begin{proof}
\par{}事实上, 上述推理把我们归结到对所有\(s\in  S\)都有不等式\hyperref[fga3.iv-3-equation-star]{公式~\ref*{fga3.iv-3-equation-star}}的情形 (必要时, 用使此不等式成立的\(s\)组成的开子集替代\(S\)). 此时可假设\(M\)纳入一个正合列\[         \mathcal {O}_S^q \to  \mathcal {O}_S^r \to  M \to  0;\] 而对于\(S\)上的\(S'\), 所要求的条件就是说在相应正合列\(\mathcal {O}_{S'}^q\to \mathcal {O}_{S'}^r\to  M'\to  0\)中, 第二个箭头为同构, 即第一个箭头为零. 于是可见, 由定义同态\(\mathcal {O}_S^q\to \mathcal {O}_S^r\)的矩阵的系数生成的理想所定义的\(S\)的闭子预概形\(Z\), 满足所求条件.
\end{proof}
\end{lemma}
\par{}回到\hyperref[fga3.iv-3-lemma-3.4]{引理~\ref*{fga3.iv-3-lemma-3.4}}的证明中刚才停下之处, 以\(Z_n\)表示由\hyperref[fga3.iv-3-lemma-3.6]{引理~\ref*{fga3.iv-3-lemma-3.6}}给模\(M_n\)和整数\(r=P(n)\)配以的\(S\)的\emph{闭}子预概形; 以\(Z'_N\)表示\(\nu \leqslant  n\leqslant  N\)时各\(Z_n\)的下确界. 则各\(Z_N\)组成\(Z\)的一列递减闭子预概形, 因而必然最终稳定. 设\(Z\)为充分大的\(N\)时各\(Z_W\)的恒定值. 这就是\hyperref[fga3.iv-3-lemma-3.4]{引理~\ref*{fga3.iv-3-lemma-3.4}}所要求的\(Z\). 至此证明了\hyperref[fga3.iv-3-lemma-3.4]{引理~\ref*{fga3.iv-3-lemma-3.4}}, 因而也证明了\hyperref[fga3.iv-3-lemma-3.3]{引理~\ref*{fga3.iv-3-lemma-3.3}}.
\par{}这样, 我们已证明\emph{\(\mathscr {Q}\kern  -.75pt\mathit {uot}_{{\mathscr {F}/X/S}}^P\)可由一个\(S\)-预概形\(Q\)表示, 后者是\(S\)上的拟射影开子预概形\(Q_\nu \)的递增并}. 为继续推进, 需要引用\hyperref[fga3.iv-2-theorem-2.1]{定理~\ref*{fga3.iv-2-theorem-2.1}}, 很容易从中推出\(Q\)是\emph{拟紧}的 (因为它是一个在\(S\)上有限型的预概形\(S'\)的像, 后者参数化各\(\mathscr {F}_K\)的Hilbert多项式为\(P\)的商层族). 因此\(Q\)等于某个\(Q_\nu \), 从而在\(S\)上拟射影. 为证明它在\(S\)上射影, 只须证明它在\(S\)上\emph{固有}; 为此只需应用[\hyperref[GD1960]{5}, II, 7.3.8]所给形式的固有性赋值判据. 只须验证:
\begin{lemma}{3.7}\setnumber{3.7}\label{fga3.iv-3-lemma-3.7}\skipsinglepar
\par{}设\(S\)是离散赋值环的谱, \(s\)是其一般点, \(X\)是\(S\)上的预概形, \(\mathscr {F}\)是\(X\)上的拟凝聚模, \(\mathscr {G}_s\)是\(X_s\)上\(\mathscr {F}_s=\mathscr {F}\otimes _{\mathcal {O}_S} k(s)\)的拟凝聚商模. 则存在唯一的\(\mathscr {F}\)的拟凝聚商模\(\mathscr {G}\), 它在\(S\)上平坦, 且其在\(X_s\)上的限制为\(\mathscr {G}_s\).
\begin{proof}
\par{}事实上, 若\(\mathscr {G}_s=\mathscr {F}_s/\mathscr {H}_s\), 只需考虑在一般纤维诱导\(\mathscr {H}_s\)的\(\mathscr {F}\)的最大子层\(\mathscr {H}\) ([\hyperref[GD1960]{5}, I, 9.4.2]), 并取\(\mathscr {G}=\mathscr {F}/\mathscr {H}\). 容易验证这个层满足要求.
\end{proof}
\end{lemma}
\par{}至此, \hyperref[fga3.iv-3-theorem-3.2]{定理~\ref*{fga3.iv-3-theorem-3.2}}以及\hyperref[fga3.iv-3-theorem-3.1]{定理~\ref*{fga3.iv-3-theorem-3.1}}均已完全证明.
\par{}这个证明同时还表明:
\begin{proposition}{3.8}\setnumber{3.8}\label{fga3.iv-3-proposition-3.8}\skipsinglepar
\par{}在\hyperref[fga3.iv-3-theorem-3.2]{定理~\ref*{fga3.iv-3-theorem-3.2}}的条件下, 置\(Q=\underline {\operatorname {Quot}}_{{\mathscr {F}/X/S}}^P\)、\(X_Q=X\times _S Q\)、\(\mathscr {F}_Q=\mathscr {F}\otimes _{\mathcal {O}_S}\mathcal {O}_Q\), 并设\(\mathscr {G}\)是\(\mathscr {F}_Q\)的一个在\(Q\)上平坦、相对Hilbert多项式为\(P\)的凝聚商, 使得\((Q,\mathscr {G})\)表示函子\(\mathscr {Q}\kern  -.75pt\mathit {uot}_{{\mathscr {F}/X/S}}^P\). 则存在一个整数\(\nu \), 使得当\(n\geqslant \nu \)时, \((f_Q)_*(\mathscr {G} (n))\)是\(Q\)上秩为\(P(n)\)的局部自由模, 并关于\(S\)非常丰沛, 即定义\(Q\)到\(S\)上某个Grassmann概形\(\underline {\operatorname {Grass}}_{P(n)} (M)\)的一个浸入. 因而当\(n\geqslant \nu \)时, \(Q\)上的层\(\bigwedge ^{P(n)} (f_Q)_*(\mathscr {G} (n))\)可逆, 并关于\(S\)非常丰沛.
\begin{proof}
\par{}事实上, 如\hyperref[fga3.iv-3-theorem-3.2]{定理~\ref*{fga3.iv-3-theorem-3.2}}中一样可归结到\(\mathscr {F}\)在\(S\)上平坦的情形, 此时只需取整数\(\nu \), 使\(A_\nu =A\) (记号同上).
\end{proof}
\end{proposition}
\par{}\hyperref[fga3.iv-3-theorem-3.2]{定理~\ref*{fga3.iv-3-theorem-3.2}}最重要的应用是\(\mathscr {F}=\mathcal {O}_X\)的情形. 此时写作\[     
\begin {aligned}       \underline {\operatorname {Quot}}_{\mathcal {O}_X/X/S}       &= \underline {\operatorname {Hilb}}_{X/S}     \\\underline {\operatorname {Quot}}_{\mathcal {O}_X/X/S}^P       &= \underline {\operatorname {Hilb}}_{X/S}^P     
\end {aligned},\] 因而有分解\[     \underline {\operatorname {Hilb}}_{X/S}     = \coprod _P \underline {\operatorname {Hilb}}_{X/S}^P.\] 按定义, \(\underline {\operatorname {Hilb}}_{X/S}\)表示函子\(\mathscr {H}\kern  -1.75pt\mathit {ilb}_{X/S} (S')\), 其值为\(X'=X\times _S S'\)的各个在\(S\)上平坦的闭子预概形的集合; \(\underline {\operatorname {Hilb}}_{X/S}^P\)则表示与具有给定Hilbert多项式\(P\)的闭子预概形相应的子函子. 这两个预概形也分别称为\(X\)在\(S\)上的\emph{Hilbert预概形}和\emph{指标为\(P\)的Hilbert预概形}. 这一术语来自Hilbert多项式在理论中所起的作用. 它们与经典Chow簇 (其目的是参数化闭链, 而非簇) 的性质差别, 类似于一个簇的闭链类的\emph{Chow环}与该簇的\emph{层的类环} (如Riemann–Roch定理\hyperref[BS1958]{[1]}中所引入的) 的差别. 注意当\(X=\mathbb {P}_S^r\)且\(S\)是一个域的谱时, 知道\(X\)上一个凝聚模\(\mathscr {F}\)的Hilbert多项式, 等价于知道\(\mathscr {F}\)的Chern类, 甚至等价于知道\(\mathscr {F}\)在\(X\)上凝聚层的类环中的类.
\begin{remarks}{3.9}\setnumber{3.9}\label{fga3.iv-3-remarks-3.9}\skipsinglepar
\par{}还应指出, \(\underline {\operatorname {Quot}}_{{\mathscr {F}/X/S}}\)和\(\underline {\operatorname {Quot}}_{{\mathscr {F}/X/S}}^P\)的构造已归结到\(X=\mathbb {P}_S^r\)且\(\mathscr {F}=\mathcal {O}_X^N\)、\(\mathcal {O}_X(1)\)为通常的非常丰沛层的情形; 更准确地说, 一般的\(\underline {\operatorname {Quot}}_{{\mathscr {F}/X/S}}\)都是上述预概形的闭子预概形. 由于各\(\underline {\operatorname {Quot}}_{{\mathscr {F}/X/S}}^P\)的构造显然与基底变换\(S'\to  S\)相容, 所以还可进一步归结到\(S=\operatorname {Spec}(\mathbb {Z})\)的情形: \[       \mathscr {Q}_{r,N}^P       = \underline {\operatorname {Quot}}_{(\mathcal {O}_{\mathbb {P}_{\mathbb {Z}}^r})^N/\mathbb {P}_{\mathbb {Z}}^r/\operatorname {Spec}(\mathbb {Z})},\] 特别地, 得到绝对Hilbert概形\[       \underline {\operatorname {Hilb}}_r^P       = \mathscr {Q}_{r,1}^P.\] 更详细地研究这些概形将很有意思, 首先是确定它们的连通分支 (它们是否连通?) 和不可约分支 (由Serre \hyperref[Ser1961]{[7]}, 可能存在完全位于某个素数\(p\neq 0\)之上的不可约分支). 回顾Weil提出的问题: \(\underline {\operatorname {Hilb}}_r^P\)在各\(s\in \operatorname {Spec}(\mathbb {Z})\)之上的纤维的不可约分支, 是否对应于素域的"正则"扩张, 即是否"相对连通". 对Hilbert概形, 这些问题或许比对"Chow簇"更容易处理. \emph{[Comp.]} Hartshorne研究了代数闭域上Hilbert概形的连通分支, 证明各\(\underline {\operatorname {Hilb}}_r^P\)连通, 并确定了使\(\underline {\operatorname {Hilb}}_r^P\neq \varnothing \)的数对\((r,P)\) \hyperref[Har1966]{[Har1966]}.
\end{remarks}
\chapter{变体}\thispagestyle{fancy}\label{fga3.iv-4}
\begin{enumerate}

  
 \item[a.]
    
\par{}\label{fga3.iv-4.a} 在\hyperref[fga3.iv-3-theorem-3.1]{定理~\ref*{fga3.iv-3-theorem-3.1}}的条件下, 设\(U\)是\(X\)的开子集; 以\(A'\)表示\(A=\mathscr {Q}\kern  -.75pt\mathit {uot}_{{\mathscr {F}/X/S}}\)的一个子函子, 使\(A'(S')\)为\(\mathscr {F}'\)的所有在\(S'\)上平坦且支集包含于\(U'\)的商模\(\mathscr {G}\)所组成的集合. 立即可见\(A'\)可由表示\(A\)的预概形\(\underline {\operatorname {Quot}}_{{\mathscr {F}/X/S}}\)的一个开子集表示. 因此, 若把\(X\)在\(S\)上射影改为在\(S\)上\emph{拟射影}, 并把结论中的"射影"一词改为"拟射影", 再以\(\mathscr {Q}\kern  -.75pt\mathit {uot}_{{\mathscr {F}/X/S}}(S')\)表示\(\mathscr {F}'\)的所有在\(S'\)上平坦且\emph{其支集在\(S'\)上固有}的凝聚商\(\mathscr {G}\)的集合, 则\hyperref[fga3.iv-3-theorem-3.1]{定理~\ref*{fga3.iv-3-theorem-3.1}}和\hyperref[fga3.iv-3-theorem-3.2]{定理~\ref*{fga3.iv-3-theorem-3.2}}仍然成立.
  

  
 \item[b.]
    
\par{}\label{fga3.iv-4.b} 一般地, 对\(\mathscr {F}'=\mathscr {F}\otimes _{\mathcal {O}_S}\mathcal {O}_{S'}\)的在\(S'\)上平坦的商\(\mathscr {G}\), 可以再加上各种在基底变换下稳定的自然条件, 从而得到我们想要表示的\(\mathscr {Q}\kern  -.75pt\mathit {uot}_{{\mathscr {F}/X/S}}\)的各种子函子. 在许多情形中, 通常的判据可以证明这些函子仍然由\(\underline {\operatorname {Quot}}_{{\mathscr {F}/X/S}}\)的开子集表示. 特别地, 加上下列某一性质时就是如此: 
\begin{enumerate}\item{}在各纤维\(X'_{s'}\)上诱导的模\(\mathscr {G}_{s'}\) (\(s'\in  S'\)) 的相伴素闭链的维数属于给定的整数集合.
        \item{}(在\(\mathscr {F}=\mathcal {O}_X\)的情形, 此时\(\mathscr {G}\)对应于\(X'\)的一个闭子预概形\(Y\):) \(Y\)是在\(S\)上\emph{光滑}的预概形 [\hyperref[Gro1960b]{3}, IV], 或者在\(S\)上\emph{正规} (即各纤维\(Y_{s'}\)在"\(k(s)\)上"正规, 也就是任意基域扩张后均正规), 或者 (若\(X\)在\(S\)上平坦) 是\(X\)中关于\(S\)的局部完全\(k\)-交 (即各纤维\(Y_{s'}\)是各\(X_{s'}\)中的局部完全交).
\end{enumerate} 其他条件可涉及在各\(X'_{s'}\)上诱导的模\(\mathscr {G}_{s'}\)的上同调性质, 等等. 当然, 如果每个条件都由\(\underline {\operatorname {Quot}}_{{\mathscr {F}/X/S}}\)的一个开集\(U_i\)表示, 那么同时满足这些条件就由这些开集之交表示. 例如, 对\(S\)上的每个\(S'\), 考虑\(X'=X\times _S S'\)的闭子预概形\(Y\)所组成的集合, 其中要求它们是在\(S'\)上具有给定秩\(r\)的平展覆盖 [\hyperref[Gro1960b]{3}, I]; 这样就得到一个关于\(S'\)的可表示反变函子.
  

  
 \item[c.]
    
\par{}\label{fga3.iv-4.c} \hyperref[fga3.ii-c.2]{[2], II, C, §2}所定义的预概形\(\underline {\operatorname {Hom}}_S(X,Y)\)、\(\prod _{X/S}Z/S\)和\(\underline {\operatorname {Isom}}_S(X,Y)\), 在适当射影性假设下存在, 并可实现为适当的Hilbert预概形的开子预概形. 因为\(\underline {\operatorname {Hom}}_S(X,Y)=\prod _{X/S}((X\times  Y)/X)\), 所以\(\operatorname {Hom}_S(X,Y)\)的情形归结到\(\prod _{X/S}(Z/X)\)的情形. 此时注意, 对\(S\)上的每个\(S'\), \(Z'=Z\times _S S'\)在\(X'=X\times _S S'\)上的截面集合, 与\(Z\)的满足下述条件的子预概形\(\Gamma \)的集合双射对应 (若\(Z\)在\(X\)上分离, 则这些子预概形必然闭): 由\(Z'\to  X'\)诱导的态射\(\Gamma \to  X'\)是同构. 因而, \emph{若\(X\)在\(S\)上平坦且固有, \(Z\)在\(S\)上拟射影, 则\(\prod _{X/S}(Z/X)\)存在, 并实现为\(\underline {\operatorname {Hilb}}_{Z/S}\)的开子预概形}. 于是, \emph{若\(X\)在\(S\)上射影且平坦, \(Y\)在\(S\)上拟射影, 则\(\underline {\operatorname {Hom}}_S(X,Y)\)存在, 并实现为\(\underline {\operatorname {Hilb}}_{(X\times _S Y)/S}\)的开子预概形}. 若\(X\)和\(Y\)都在\(S\)上射影, 则立即得到\(\underline {\operatorname {Isom}}_S(X,Y)\)也存在, 并由\(\underline {\operatorname {Hom}}_S(X,Y)\)的一个开子集表示. 同样, 若\(X\)在\(S\)上平坦且射影, \(Y\)在\(S\)上拟射影, 则\(S\)-预概形\(\underline {\operatorname {Imm}}_S(X/Y)\)也可由\(\underline {\operatorname {Hom}}_S(X,Y)\)的开子集表示; 这个预概形对应于\(\underline {\operatorname {Hom}}_S(X,Y)\)所表示的函子的一个子函子, 后者对应于那些为浸入的\(S'\)-态射\(X'\to  Y'\).

    
\par{}设\(\mathscr {L}\) (分别为\(\mathscr {M}\)) 是\(X\) (分别为\(Y\)) 上关于\(S\)非常丰沛的可逆层, 从而在\(X\times _S Y\)上得到关于\(S\)非常丰沛的层\(\mathscr {L}\otimes _{\mathcal {O}_S}\mathscr {M}\). 对任意有理系数多项式\(P\), \(\underline {\operatorname {Hilb}}_{(X\times _S Y)/S}^P\)有定义, 且为\(S\)上的拟射影预概形. 它因而在\(\underline {\operatorname {Hom}}_S(X,Y)\)上诱导一个既开又闭、并在\(S\)上拟射影的子集, 记作\(\operatorname {Hom}_S(X,Y)^P\). 于是\(\underline {\operatorname {Hom}}_S(X,Y)^P\)在\(S\)上的截面, 就是那些使任意整数\(n\)均满足\[         \chi \Big ((\mathscr {L}\otimes _{\mathcal {O}_X}g^*(\mathscr {M}))^{\otimes  n}\Big )         = P(n).       \]的\(S\)-态射\(g\colon  X\to  Y\). 这样得到Matsusaka定理的推广: 一个"极化"射影簇的自同构组成一个代数群. 这个断言在这里显然具有更精确的意义, 因为该群被定义为一个泛问题的解. 还应注意, 在代数闭域上, 过去所考虑的自同构群, 是将这里定义的"真正"自同构群模掉幂零元所得的群; 这说明为何那些历史上的构造难以在非完美基域上进行, 因为基域扩张后出现的幂零元理想未必"在\(k\)上有定义". 同样的说明也适用于过去的大多数构造.
  

\end{enumerate}
\chapter{Hilbert概形的微分研究}\thispagestyle{fancy}\label{fga3.iv-5}
\par{}有以下结果:
\begin{proposition}{5.1}\setnumber{5.1}\label{fga3.iv-5-proposition-5.1}\skipsinglepar
\par{}设\(S\)是预概形, \(S_0\)是由平方为零的拟凝聚理想\(\mathscr {I}\)定义的子预概形, \(X\)是\(S\)-预概形, \(\mathscr {F}\)是\(X\)上的拟凝聚模. 置\(X_0=\mathscr {F}\times _S S_0\)及\(\mathscr {F}_0=\mathscr {F}\otimes _{\mathcal {O}_S}\mathcal {O}_{S_0}\). 最后, 设\(\mathscr {G}_0=\mathscr {F}_0/\mathscr {H}_0\)是\(\mathscr {F}_0\)的一个在\(S_0\)上平坦的拟凝聚商模. 对\(X\)的每个开子集\(U\), 以\(\mathscr {E}(U)\)表示\(\mathscr {F}|U\)的所有在\(S\)上平坦并满足\(\mathscr {G}\otimes _{\mathcal {O}_S}\mathcal {O}_{S_0}=\mathscr {G}_0\)的拟凝聚商模\(\mathscr {G}\)所组成的集合; 当\(U\)变动时, 各\(\mathscr {E}(U)\)是\(U\)上一个层\(\mathscr {E}\)的截面. 此时群层\[       \mathscr {A}       = \underline {\operatorname {Hom}}_{\mathcal {O}_{X_0}}(\mathscr {H}_0,\mathscr {G}_0\otimes _{\mathcal {O}_{S_0}}\mathscr {I})     \]自然作用于\(\mathscr {E}\), 使它成为一个"形式地为\(\mathscr {A}\)-主齐性"的层 (即对于\(X\)的每个开子集\(U\), \(\mathscr {E}(U)\)或者为空, 或者是一个\(\mathscr {A}(U)\)-主齐性集).
\end{proposition}
\par{}因此得到:
\begin{corollary}{5.2}\setnumber{5.2}\label{fga3.iv-5-corollary-5.2}\skipsinglepar
\par{}假设在\(X\)上局部存在将\(\mathscr {G}_0\)延拓为\(\mathscr {F}\)的在\(S\)上平坦的商的延拓\(\mathscr {G}\) (即层\(\mathscr {E}\)的各茎非空). 则存在一个典范的障碍类\[       c(\mathscr {G}_0) \in  \operatorname {H}^1(X,\mathscr {A}),\] 它为零的充要条件是存在把\(\mathscr {G}_0\)延拓为\(\mathscr {F}\)的在\(S\)上平坦的商的整体延拓\(\mathscr {G}\). 若这个类为零, 则所有可能延拓组成的集合\(\mathscr {E}(X)\)是\(\mathscr {A}(X)=\operatorname {Hom}_{\mathcal {O}_X}(\mathscr {H}_0,\mathscr {G}_0\otimes _{\mathcal {O}_{S_0}}\mathscr {I})\)下的一个主齐性集.
\end{corollary}
\par{}特别地, 若\(\operatorname {H}^1(X,\mathscr {A})=0\), 就可保证整体延拓存在.
\begin{corollary}{5.3}\setnumber{5.3}\label{fga3.iv-5-corollary-5.3}\skipsinglepar
\par{}假设\(Q=\underline {\operatorname {Quot}}_{\mathscr {F}/X/S}\)存在 (参见\hyperref[fga3.iv-4.a]{§4.a}); 例如可假设\(X\)在某个局部Noether的\(S\)上拟射影, 且\(\mathscr {F}\)凝聚. 设\(x\in  Q\), 它对应于某个\(k(s)\)的剩余域扩张\(K=k(x)\) (其中\(s\in  S\)). 则\(x\)由\(K\)-预概形\(X_K\)上的模\(\mathscr {F}_0=F_K\)的一个凝聚商模\(\mathscr {G}_0=\mathscr {F}_0/\mathscr {H}_0\)定义. 设\(\mathscr {A}\)是\(X_K\)上由\[       \mathscr {A} = \underline {\operatorname {Hom}}_{\mathcal {O}_{X_0}}(\mathscr {H}_0,\mathscr {G}_0).     \]定义的凝聚层. 则纤维\(Q_s\)在点\(x\)处的Zariski切空间 (即\(\mathfrak {m}/\mathfrak {m}^2\)的\(K\)-对偶, 其中\(\mathfrak {m}\)是\(\mathcal {O}_{Q_K,x}\)的极大理想) 典范同构于\(\operatorname {H}^0(X_K,\mathscr {A})\).
\end{corollary}
\par{}这个关于Zariski切空间的结果可以推广: 对给定的\(S\)-态射\(g\colon  S'\to  Q\), 即\(Q'=Q\times _S S'\)在\(S'\)上的一个截面\(g'\), 可用以下公式刻画模\[     \Omega      = g^*(\Omega _{Q/S}^1)     = {g'}^*(\mathscr {J}/\mathscr {J}^2)   \] (其中\(\mathscr {J}\)是由\(Q'\)在\(S'\)上的截面\(g'\)所定义的\(Q'\)上的理想): \[     \operatorname {Hom}_{\mathcal {O}_{S'}}(\Omega ,\mathscr {M})     \simeq  \operatorname {H}^0(X',\mathscr {A});\] 这个公式关于\(S'\)上的凝聚模\(\mathscr {M}\)函子化, 其中\(\mathscr {A}\)仍为\(X'=X\times _S S'\)上由\[     \mathscr {A} = \underline {\operatorname {Hom}}_{\mathcal {O}_{X'}}(\mathscr {H},\mathscr {G}\otimes _{\mathcal {O}_S}\mathscr {M})   \]定义的模 (\(\mathscr {G}=\mathscr {F}'/\mathscr {H}\)是\(\mathscr {F}'=\mathscr {F}\otimes _{\mathcal {O}_S}\mathcal {O}_{S'}\)中与\(g\)相应的商模). 只需在\hyperref[fga3.iv-5-proposition-5.1]{命题~\ref*{fga3.iv-5-proposition-5.1}}中以\(S'\)替代\(S_0\), 并以预概形\(D(\mathscr {M})=(S',\mathcal {O}_{S'}+\mathscr {M})\)替代\(S\), 其中把\(\mathscr {M}\)视为平方为零的理想.
\par{}若在\hyperref[fga3.iv-5-proposition-5.1]{命题~\ref*{fga3.iv-5-proposition-5.1}}中\(\mathscr {F}=\mathcal {O}_X\), 则给定\(\mathscr {G}_0\)等价于给定\(X_0\)的一个在\(S_0\)上平坦的闭子预概形\(Y_0\), 其定义理想为\(\mathscr {J}_0=\mathscr {M}_0\), 而\((*)\)给出\[     \mathscr {A}     = \underline {\operatorname {Hom}}_{\mathcal {O}_{X_0}}(\mathscr {J}_0/\mathscr {J}_0^2,\mathcal {O}_{Y_0}\otimes _{\mathcal {O}_{S_0}}\mathscr {J}).\] 此时把\(\mathscr {J}/\mathscr {J}^2\)视为\(Y_0\)在\(X_0\)中的\emph{余法层}, 亦记作\(\mathscr {N}_{Y_0/X_0}\); 将\(\mathscr {A}\)看成\(Y_0\)上的模, 并在\(Y\)上计算\(\operatorname {H}^0\)和\(\operatorname {H}^1\), 是有利的. 若\(Y_0\)局部是\(X_0\)中的完全交, 且\(X\)在\(S\)上平坦, 则在\hyperref[fga3.iv-5-proposition-5.1]{命题~\ref*{fga3.iv-5-proposition-5.1}}中局部延拓的可能性可得到保证, 且\(\mathscr {J}/\mathscr {J}^2\)在\(Y_0\)上局部自由, 可以写作\[     \mathscr {A}     = \check {\mathscr {N}}_{X_0/Y_0}\otimes _{\mathcal {O}_{S_0}}\mathscr {J};\] 其中右侧第一因子是\(Y_0\)在\(X_0\)中的法丛. 应用光滑性的基本判据 [\hyperref[Gro1960b]{3}, III, 3.1], 例如得到:
\begin{corollary}{5.4}\setnumber{5.4}\label{fga3.iv-5-corollary-5.4}\skipsinglepar
\par{}在\hyperref[fga3.iv-5-corollary-5.3]{推论~\ref*{fga3.iv-5-corollary-5.3}}的条件下, 假设\(\mathscr {F}=\mathcal {O}_X\), \(X\)在\(S\)上平坦, 且与\(\mathscr {G}_0\)相应的\(X_0\)的闭子预概形\(Y_0\)局部为完全交. 则\(Q_s\)在\(x\)处的Zariski切空间典范同构于\(\operatorname {H}^0(Y_0,\check {\mathscr {N}}_{X_0/Y_0})\). 若\(\operatorname {H}^1(Y_0,\check {\mathscr {N}_{X_0/Y_0}})=0\), 则Hilbert预概形\(X\)在点\(x\)处关于\(S\)光滑 (其中\(\check {\mathscr {N}}_{X_0/Y_0}\)是\(Y_0\)在\(X_0\)中的法层).
\end{corollary}
\begin{remark}{5.5}\setnumber{5.5}\label{fga3.iv-5-remark-5.5}\skipsinglepar
\par{}这个断言特别适用于\(Y_0\)是\(X_0\)中由\emph{一}个方程定义的完全交的情形, 即它是一个正"Cartier除子". 此时\(\check {\mathscr {N}}_{X_0/Y_0}\)同构于由除子\(Y_0\)所定义的\(X_0\)上的可逆层\(\mathscr {J}^{-1}\)在\(Y_0\)上诱导的层. 在研究非奇异射影簇\(X_0\)上的正除子族时, 特别会遇到这种情形. \(Q\) (或若愿意, \(Q\)中对应于除子的开集\(D\)) 在点\(x\)处的Zariski切空间与\(\mathscr {H}^0(Y_0,\check {\mathscr {N}}_{X_0/Y_0})\)的同构, 在经典代数几何中称为"特征同态" (从前者到后者). 过去只在\(x\)是某个除子的"完全连续族"的参数簇\(T\)的光滑点时定义它; 从我们的观点看, 此参数簇就是概形\(D\)的某个不可约分支, 赋以诱导的\emph{约化}结构. \(T\)在\(x\)处的切空间于是为\(D\)在\(x\)处切空间的一个\emph{子空间}, 所以旧式特征同态确实单射, 但只有加上额外条件才满射, 例如\(D\)在\(x\)处整时. 事实上, Zappa \hyperref[Zap1945]{[8]}构造过一个例子 (其中\(X\)是\emph{复数域}上的非奇异射影曲面), 使特征同态即使在\(T\)的一般点处也不满射. \emph{因此\(D\)即使在所讨论不可约分支的一般点处也不是整的} 这格外鲜明地表明, 为理解曲面理论中最经典的现象, 必须容许带幂零元的簇.
\par{}\emph{[Comp.]} 关于Zappa的例子, 注意Mumford甚至已构造\(\mathbb {P}_\mathbb {C}^3\)的Hilbert概形的一个不可约分支 (其一般点表示次数\(14\)、亏格\(24\)的非奇异曲线), 它在其一般点处不约化. 沿所得曲线作爆破, 他还得到一个\(\mathbb {C}\)上的\(3\)维正则射影概形, 其形式模概形在一般点处不约化; 等价地, 它在\(\mathbb {C}\)上按解析几何意义定义的局部模簇 (见\emph{Séminaire Cartan} \textbf{13}, 1960/61) 在所有点处均不约化.
\end{remark}
\begin{remark}{5.6}\setnumber{5.6}\label{fga3.iv-5-remark-5.6}\skipsinglepar
\par{}我们在\hyperref[fga3.iv-5-corollary-5.4]{推论~\ref*{fga3.iv-5-corollary-5.4}}中给出了光滑性判据, 特别适用于除子概形. Kodaira在\hyperref[Kod1956]{[6]}中给出另一判据, 即要求\(\operatorname {H}^1(X_0,\mathscr {L})\)为零, 其中\(\mathscr {L}=\mathscr {J}_0^{-1}\)是由除子\(Y_0\)定义的\(X_0\)上的可逆层; 这一判据适用于\(S\)是特征\(0\)的域的谱的情形, 而\hyperref[Kod1956]{[6]}在基域为\(\mathbb {C}\)时用超越方法证明了它. 这里指出, 一般地, 对任意\(S\), Kodaira的条件足以保证从除子预概形到\(X/S\)的Picard预概形的典范态射\(D\to \underline {\operatorname {Pic}}_{X/S}\), 在所讨论的点\(x\)处光滑 (一旦知道\(\underline {\operatorname {Pic}}_{X/S}\)存在, 用通常的光滑性判据就容易验证). 此外, 若\(\underline {\operatorname {Pic}}_{X/S}\)在\(x\)的像点处关于\(S\)光滑 (例如\(\underline {\operatorname {Pic}}_{X/S}\)在\(S\)上光滑), 则\(D\)在\(x\)处关于\(S\)光滑. 另一方面, Cartier已证明在\emph{特征\(0\)的域\(k\)}上局部有限型的每个群预概形都在\(k\)上光滑. 合用这两个结果, 就恢复了Kodaira的结果. 由这些说明还可见, 在特征\(p>0\)的域\(K\)上, 若\(\underline {\operatorname {Pic}}_{X/S}\)不在\(k\)上光滑 (例如\(X\)是Igusa曲面时), 则条件\(\operatorname {H}^1(X_0,\mathscr {L})=0\)反而蕴涵\(D\)在\(x\)处不光滑; 若\(K\)代数闭, 则它甚至在\(x\)处不约化.
\end{remark}
\par{}最后给出下面的结果, 它在纤维空间的微分研究中起重要作用:
\begin{proposition}{5.7}\setnumber{5.7}\label{fga3.iv-5-proposition-5.7}\skipsinglepar
\par{}设\(S\)局部 Noether，\(X\)是在\(S\)上有限且平坦的预概形, \(Z\)是\(S\)上的预概形, 并使\(\prod _{X/S}(Z/X)\)存在 (若\(Z\)在\(X\)上拟射影就如此). 若\(Z\)在\(X\)上光滑, 则\(\prod _{X/S}(Z/X)\)在\(S\)上光滑.
\begin{proof}
\par{}这由定义以及通常的光滑性判据 [\hyperref[Gro1960b]{3}, III, §3.1] 立即得到.
\end{proof}
\end{proposition}
\par{}注意, 若\(X\)在\(S\)上有限且平坦, 则\(\prod _{X/S}(Z/X)\)的存在性问题可用很初等的方式处理, 无须使用Hilbert概形理论. 例如可得: 若\(X\)在\(S\)上为根态射, 则无须对\(Z\)作任何限制, \(\prod _{X/S}(Z/X)\)就存在. 举例说, 设\(T\)是\(S\)-预概形, \(T_n\)是\(T\times _S T\)中对角线\(T\times _S T\)的"\(n\)阶无穷小邻域", 带有由两个投影诱导的态射\(p_1,p_2\colon  T_n\to  T\). 借助\(p_1\), 可把\(T_n\)视为\(T\)上的有限预概形; 并进一步假设\(T_n\)在\(T\)上平坦 (当\(T\)在\(S\)上光滑时即如此). 对每个\(T\)上的预概形\(X\), 置\[     (X/T/S)^{(n)}     = \prod _{T_n/S}(p_2^*(X/T)/T_n);\] 这是\(T\)上的一个预概形, 称为\emph{\(X\)在\(T\)上的\(n\)阶截面芽丛} (相对于\(S\)). 它函子地依赖于\(X\), 且若\(X\)光滑, 则它也在\(T\)上光滑.
\chapter{与范数及对称积概念的关系}\thispagestyle{fancy}\label{fga3.iv-6}
\par{}设\(S\)是预概形, \(X\)、\(Y\)是\(S\)-预概形, \[     u\colon (X/S)^n\to  Y   \]是从\(X/S\)的\(n\)重笛卡尔幂到\(Y\)的一个\emph{对称}的\(S\)-态射. 为简单起见, 假设\(S\)局部Noether, \(X\)和\(Y\)在\(S\)上有限型. 此时对于\(X\)上每个支集在\(S\)上有限、并且在\(S\)上平坦且相对\(S\)的秩等于\(n\)的凝聚模\(\mathscr {F}\) (即\(f_*(\mathscr {F})\)是\(S\)上秩为\(n\)的局部自由模), 可以自然地配以\(Y\)在\(S\)上的一个截面: \[     \mathscr {N}_{X/S}^u(\mathscr {F}) \in  \Gamma (Y/S).\] 这里不详述定义, 只指出所得的形式运算是通常范数及迹的形式运算的自然推广. 当\(X\)在\(S\)上的\(n\)重对称幂存在时 (例如对称群\(\sigma _n\)作用于\((X/S)^n\)的各轨道均包含于仿射开集时), 可取\(Y\)为这个对称幂\(\operatorname {Symm}_S^n(X)\), 从而得到一个典范元素\[     \mathscr {N}_{X/S}(\mathscr {F}) \in  \Gamma (\operatorname {Symm}_S^n(X)/S),\] 并可由它恢复各\(\mathscr {N}_{X/S}^u(\mathscr {F})\). 另一个重要情形是\(X\)为\(S\)上的交换幺半群, \(X=Y\), 且态射\(u\)来自\(X\)的合成律. 此时对于\(X\)上的模\(\mathscr {F}\)所对应的\(X\)在\(S\)上的截面, 简记为\(\mathscr {N}(\mathscr {F})\).
\par{}现在假设\(X\)上给定一个凝聚模\(\mathscr {F}\), 使\(\underline {\operatorname {Quot}}_{{\mathscr {F}/X/S}}\)存在; 或者至少使下述函子\(\mathscr {Q}\kern  -.75pt\mathit {uot}_{{\mathscr {F}/X/S}}^n\)可由一个\(S\)-预概形\(\underline {\operatorname {Quot}}_{{\mathscr {F}/X/S}}^n\)表示: 它给\(S\)上的每个\(S'\)配以\(\mathscr {F}'=\mathscr {F}\otimes _{\mathcal {O}_S}\mathcal {O}_{S'}\)的所有在\(S\)上平坦、相对秩为\(n\)的凝聚商层\(\mathscr {M}\)所组成的集合. (若\(X\)在\(S\)上拟射影, 则\(\underline {\operatorname {Quot}}_{{\mathscr {F}/X/S}}^n\)确实存在; 按\cref{fga3.iv-3}的记号, 它恰为\(\underline {\operatorname {Quot}}_{{\mathscr {F}/X/S}}^P\), 其中\(P\)是常数项为\(n\)的常值多项式.) 由于各\(\mathscr {N}_{X/S}^u(\mathscr {M})\)的构造与基底变换相容, 就得到一个典范态射\[     \mathscr {N}_{X/S}^u\colon  \underline {\operatorname {Quot}}_{{\mathscr {F}/X/S}}^n \to  Y;\] 特别地, 若\(X\)在\(S\)上的\(n\)重对称幂存在, 则有\[     \mathscr {N}_{X/S}\colon  \underline {\operatorname {Quot}}_{{\mathscr {F}/X/S}}^n \to  \operatorname {Symm}_S^n(X).\] 最重要的情形是\(\mathscr {F}=\mathcal {O}_X\), 此时得到态射\[     \mathscr {N}_{X/S}\colon  \underline {\operatorname {Hilb}}_{X/S}^n \to  \operatorname {Symm}_S^n(X).\] 显然在\(n=0\)和\(n=1\)时, 它为同构. 但对\(n\geqslant 1\), 即使\(S\)是域\(k\)的谱, 且\(X\)在\(S\)上光滑, 它一般也不是同构, 甚至不是单射态射; 因为仅知道\(X\)的\(0\)维子概形所定义的闭链, 并不能确定该子概形 (例如它对应于\(X\)的闭点\(x\)处局部环\(\mathcal {O}_{X,x}\)中关于极大理想的准素理想\(\mathscr {I}\); 准确地说, 只知道\(\mathscr {I}\)在\(\mathcal {O}_{X,x}\)中作为\(k\)-向量空间的余维). 只能作如下断言 (这里\(S\)又是任意的): 
\begin{enumerate}
    
 \item[a.]若\(X\)在\(S\)上光滑, 则范数态射在\(\underline {\operatorname {Hilb}}_{X/S}^n\)中对应于\(X\)内秩为\(n\)的平展覆盖的分类的开集 (参见\hyperref[fga3.iv-4.b]{§4.b}) 与\(\operatorname {Symm}_S^n(X)\)中对应于无重分支的\(n\)-闭链的开集之间定义一个同构.

    
 \item[b.]若此外\(X\)在\(S\)上的相对维数为\(1\), 则范数态射甚至定义从\(\underline {\operatorname {Hilb}}_{X/S}^n\)到\(\operatorname {Symm}_{X/S}^n\)的一个同构.

  
\end{enumerate} 第二点来自以下事实: 非奇异代数曲线的\(0\)维子概形, 只要知道其相应的闭链就能确定. 同样的说明更一般地适用于非奇异代数概形上的正Cartier除子 (而在这个很特殊的情形, 也不排除Chow簇与Hilbert簇给出相同结果的可能).
\chapter{补充与问题}\thispagestyle{fancy}\label{fga3.iv-7}
\par{}如J.-P. Serre所指出, 由Nagata的一个熟知例子可知, 可以找到一个为域\(k\)的谱的概形\(S\)、一个为\(k\)的二次扩张\(k'\)的谱的\(S\)-概形\(S'\), 以及一个\(3\)维的光滑固有 (但非射影) \(S'\)-概形\(X\), 使\(\prod _{S'/S}(X/S)\)不存在. 这当然更蕴涵Hilbert概形\(\underline {\operatorname {Hilb}}_{X/S}^2\)不存在 (甚至表示\(X\)内\(S\)的秩为\(2\)的平展覆盖的\(k\)-概形也不存在; 从而\(X\)的对称平方也不存在, 参见\cref{fga3.iv-6}). 这对代数几何中进行非射影构造的可能性构成了严重限制. (不过可以相信, 解析几何中不会出现这样的限制, 正如形式几何中不出现一样, 参见\hyperref[fga3.ii]{[2, II]}.) 然而, 若\(X\)是域\(k\)的谱\(S\)上的固有概形, 且\(Z\)在\(X\)上拟射影, 则\(\prod _{X/S}(Z/X)\)存在, 并且是由\(S\)上一列拟射影概形的和构成的概形 (如同射影情形\hyperref[fga3.iv-3-theorem-3.1]{定理~\ref*{fga3.iv-3-theorem-3.1}}). 为证明这一点, 可以取一个射影\(S\)-概形\(X'\)支配\(X\), 从而归结到\(X\)本身射影的情形; 这里不详述证明, 它还使用\hyperref[fga3.i-a.2.b]{[2], I, A, §2.b}中有限态射的分解结果. 这一方法成功的关键在于, 当\(S\)是域的谱时, Chow引理中出现的\(X'\)自动在\(S\)上平坦. 若对\(S\)不加假设, 而仅假设\(X\)在\(S\)上固有平坦, \(Z\)在\(X\)上拟射影, 我不知道结论是否仍成立. 应用中一个重要情形是\(Z\)为\(X\)的闭子概形; 此时若\(\prod _{X/S}(Z/X)\)存在, 它必为\(S\)的闭子概形. 只要\(X\)在\(S\)上射影, 就可以用相当简单的办法直接构造它, 无须用Hilbert概形理论; 所用方法还更一般地表明, 若\(Z\)在\(X\)上仿射, 则\(\prod _{X/S}(Z/X)\)存在且在\(S\)上仿射. 它也表明, 若\(X\)在\(S\)上固有平坦 (但不一定在\(S\)上射影), 则对\(X\)上每个局部平凡向量丛\(Z\), \(\prod _{X/S}(Z/X)\)存在且为\(S\)上的向量丛. 希望能重新研究并统一这些结果.

\section*{参考文献}

\begin{enumerate}
\item[{[1]}]\setnumber{1}\label{fga3.iv-bib-1} BOREL (A.) 与 SERRE (J.-P.). --- Le théorème de Riemann--Roch, \emph{Bull. Soc. math. France}, t. 86, 1958, p. 97--136.
\item[{[2]}]\setnumber{2}\label{fga3.iv-bib-2} GROTHENDIECK (Alexander). --- Technique de construction et théorèmes d'existence en géométrie algébrique, I, II, III, \emph{Séminaire Bourbaki}, t. 13, 1959/60, n\textsuperscript{o} 190, 29页, n\textsuperscript{o} 195, 22页; t. 14, 1960/61, n\textsuperscript{o} 212, 20页.
\item[{[3]}]\setnumber{3}\label{fga3.iv-bib-3} GROTHENDIECK (Alexander). --- \emph{Séminaire de Géométrie algébrique}, I, II, III, IV. Paris, Institut des hautes Études scientifiques, 1960/61 (油印本).
\item[{[4]}]\setnumber{4}\label{fga3.iv-bib-4} GROTHENDIECK (Alexander). --- Technique de construction en géométrie analytique, IV, V, \emph{Séminaire Cartan}, t. 13, 1960/61, n\textsuperscript{o} 11 与 12.
\item[{[5]}]\setnumber{5}\label{fga3.iv-bib-5} GROTHENDIECK (A.) 与 DIEUDONNÉ (J.). --- \emph{Éléments de géométrie algébrique}, I: Le langage des schémas. Paris, Presses universitaires de France, 1960 (Institut des hautes Études scientifiques, Publications mathématiques, 4); II, III, IV (待出版).
\item[{[6]}]\setnumber{6}\label{fga3.iv-bib-6} KODAIRA (K.). --- Characteristic linear systems of complete continuous systems, \emph{Amer. J. of Math.}, t. 78, 1956, p. 716--744.
\item[{[7]}]\setnumber{7}\label{fga3.iv-bib-7} SERRE (Jean-Pierre). --- Exemples de variétés projectives en caractéristique \(p\) non relevables en caractéristique \(0\), \emph{Proc. Nat. Acad. Sc. U. S. A.}, t. 47, 1961, p. 108--109.
\item[{[8]}]\setnumber{8}\label{fga3.iv-bib-8} ZAPPA (G.). --- Sull'esistenza, sopra la superficie algebriche, di sistemi continui completi infiniti, la cui curva e a serie caratteristica incompleta, \emph{Pont. Acad. Sc. Acta}, t. 9, 1945, p. 91--93.

\end{enumerate}

\section*{增补}
\label{fga3.iv-additif}

\begin{center}[校样校正时增补]
\end{center}

\par{}现在看来, 文献[2]的III, §8中的各个猜想都是错误的, 即使对于特征\(0\)的域上的非奇异簇, 无论商的存在性还是商的拟射影性, 都是如此; 即使该群的作用具有闭图, 也仍然如此.
